\documentclass[10pt,a4paper]{amsart}
\usepackage[margin=19mm]{geometry}
\usepackage{amsmath,amssymb,mathtools}
\usepackage[scr=rsfs,cal=euler]{mathalfa}
\usepackage{xfrac}
\usepackage[unicode,hidelinks,bookmarks=false]{hyperref}
\newcommand{\ie}{i.e.\ }
\newcommand{\cf}{cf.\ }
\newcommand{\resp}{respectively\ }
\newcommand{\Subsection}{\subsection}
\providecommand{\nobreakdash}{\nobreak\hbox{-}}
\setlength{\emergencystretch}{2em}
\multlinegap0pt
\usepackage{amssymb}
\usepackage{xcolor}
\usepackage{mathtools}
\usepackage{tikz}
\usepackage{cancel}
\usepackage{float}
\usepackage{bm}
\newtheorem{cor}{Corollary}
\newtheorem{prop}{Proposition}
\newcommand{\C}{\ensuremath{\mathbb C}}
\newcommand{\CP}{\ensuremath{\mathbb {CP}}}
\newcommand{\cE}{{\mathcal E}}
\newcommand{\cO}{\ensuremath{\mathcal O}}
\newcommand{\de}{\mathrm{d}}
  \def\mathbf#1{#1}%
\title{ALG gravitational instantons and Hitchin moduli spaces, I: \\ Torelli parameters}
\author{Laura Fredrickson, Rafe Mazzeo, Jan Swoboda,, Hartmut Weiss}
\date{Source: arXiv:2603.17020v1 (CC BY 4.0)}
\begin{document}
\maketitle

\noindent\emph{Source and licence.} ALG gravitational instantons and Hitchin moduli spaces, I: \\ Torelli parameters. Version arXiv:2603.17020v1 (CC BY 4.0), distributed by arXiv under the Creative Commons Attribution 4.0 International licence, http://creativecommons.org/licenses/by/4.0/, as recorded in the arXiv licence metadata for this exact version and shown on the abstract page https://arxiv.org/abs/2603.17020v1. Full licence notice: \texttt{licenses/arxiv-2603.17020-CC-BY-4.0.txt}.\par\smallskip

\noindent\textit{Editorial scope.} Counted unit: the article's prop prop:interiorintersection (source0.tex lines 3072--3082). The article's complete original proof of it is delivered with it (source0.tex lines 3084--3213, one \texttt{proof} environment of 130 lines). No mathematics is written, repaired or re-derived: the statement and the proof are the article's own lines, in the article's own notation, and no retained line is substituted or re-flowed.  Retained uncounted as the source-local context the proof needs: lines 1319--1319; lines 1521--1521; lines 2696--2699; lines 2703--2711; lines 2715--2718; lines 2729--2754; lines 2962--2967.  Nothing was omitted from any retained span, so the package carries no omission record.  No retained line contains a citation, so the wrapper carries no bibliography and no bibliography entry.  No retained line contains a figure environment, an image-inclusion command or drawing code, so no image is delivered and the package declares no asset dependency.  Editorial formatting only, in addition: an AMS \texttt{amsart} preamble that restates the article's own declarations and macros used by the retained lines, namely the environments \texttt{cor}, \texttt{prop} on separate counters, together with the standard packages those lines need.  The retained environments print in this document as Corollary 1, Proposition 1 in order of appearance, whereas the article prints the same items with the numbering of its own full document.  The complete native source of the article as distributed in the arXiv e-print for this exact version is bundled unchanged as \texttt{sources/196596/source0.tex}.
\subsubsection{Description of exterior chambers}\label{sec:exteriorchambers} The combinatorial labels for the exterior chambers are the subsets 

\subsubsection{Central Sphere in Interior Chambers} \label{sec:centralsphereinint}

 \begin{align}\label{eq:bigstrata}
 \cE &\simeq \cO(-2) \oplus \cO(-2)  \\ \nonumber
 \varphi_{\beta, u, w} &= \frac{\de z}{z(z-1)(z-p_0)}\begin{pmatrix}  -(m_\infty z^2 + w) & \frac{f(z) + \beta z (z-1)(z-p_0) - (m_\infty z^2 + w)^2}{z-u} \\  z-u & m_\infty z^2 + w \end{pmatrix},
\end{align}

\begin{align} \label{eq:extrapoint}
  \cE &\simeq \cO(-2) \oplus \cO(-2)  \\ \nonumber
 \varphi_\beta &=  \frac{\de z}{z(z-1)(z-p_0)}\begin{pmatrix} -(m_\infty z^2 +
            \frac{f_3 + \beta}{2 m_\infty} z) & f(z) + \beta z(z-1)(z-p_0) 
            -\left(m_\infty z^2 +
            \frac{f_3 + \beta}{2 m_\infty} z\right)^2 \\ 1 & m_\infty z^2 +
            \frac{f_3 + \beta}{2 m_\infty} z.
           \end{pmatrix}.
\end{align}

\begin{align}\label{eq:smallstrata}
 \cE \simeq&  \cO(-1) \oplus \cO(-3) \\ \nonumber
 \varphi_\beta  =& \frac{\de z}{z(z-1)(z-p_0)}\begin{pmatrix} 0 & f(z) + \beta z (z-1)(z-p_0) \\ 1 & 0 \end{pmatrix} 
\end{align}

\begin{cor} [Flags] \label{cor:flags}
For the Higgs bundles in the big strata \eqref{eq:bigstrata}, the flags are 
     \begin{equation}\label{eq:flags}
    F_0 = \begin{pmatrix} \frac{-m_0 p_0 +w}{u} \\ 1 \end{pmatrix}, \quad 
      F_1 = 
  \begin{pmatrix} \frac{m_1(p_0-1) + m_\infty +w}{u-1} \\ 1 \end{pmatrix}, \quad 
        F_{p_0} = \begin{pmatrix} \frac{-m_{p_0} p_0(p_0-1) + m_\infty p_0^2 +w}{u-p_0}\\ 1 \end{pmatrix}, \quad 
          F_\infty = \begin{pmatrix} 1\\ 0 \end{pmatrix};
   \end{equation}
 for the Higgs bundles in the additional point of the big stratum \eqref{eq:extrapoint} these are   \begin{align*}
    F_0 &= \begin{pmatrix} m_0 p_0  \\ 1 \end{pmatrix}, \quad 
      F_1 = 
  \begin{pmatrix} - \frac{f_3 + 2 m_\infty (m_\infty + m_1(p_0-1)) + \beta}{2 m_\infty} \\ 1 \end{pmatrix}, 
        F_{p_0} &= \begin{pmatrix} - \frac{f_3 p_0+ 2 m_\infty m_{p_0}(p_0 - 1) p_0 + p_0(2 m_\infty^2 p_0 + \beta)}{2m_\infty}\\ 1 \end{pmatrix}, \quad 
          F_\infty = \begin{pmatrix} 1\\ 0 \end{pmatrix},
   \end{align*}
   where $f_3$ is shorthand for the coefficient of $z^3$ in $f_{\mathbf{m}}(z)$;
 and for the Higgs bundles in the small stratum \eqref{eq:smallstrata}
  \begin{equation*}
    F_0 = \begin{pmatrix} m_0 p_0 \\ 1 \end{pmatrix}, \quad 
      F_1 = 
  \begin{pmatrix} m_1(1-p_0)  \\ 1 \end{pmatrix}, \quad 
        F_{p_0} = \begin{pmatrix} m_{p_0} p_0(p_0-1) \\ 1 \end{pmatrix}, \quad 
          F_\infty = \begin{pmatrix} - m_\infty\\ 1 \end{pmatrix}.
   \end{equation*}
\end{cor}

\begin{prop}\label{prop:exteriorintersection}
In the exterior chamber labelled by $D_{I_0} = \{0, 1, p_0\}$ (see Section \ref{sec:exteriorchambers}) 
  \[\int_{\CP^1_p} \Omega_I = \begin{cases} -4 \pi  m_p &\quad p \in \{0, 1, p_0\}\\
  4 \pi  m_p & \quad p=\infty\end{cases},  \]
  i.e. $\int_{\CP^1_J} \Omega_I = 2 \pi (M_J -M_{I_0}).$
 \end{prop}

\begin{prop}\label{prop:interiorintersection} 
In the interior chamber of Type B2 with $\mathcal{I}=\{\{0, 1, p_0, \infty\}, \{0, 1\}, \{0, p_0\}, \{1, p_0\}\}$ then we have  
\begin{equation}
\begin{array}{llll}
\int_{S^2_\infty} \Omega_I &=\int_{S^2_{\{0, 1, p_0, \infty \}}} \Omega_I &= -2 \pi (-m_0 - m_1 - m_{p_0}- m_\infty)&=2 \pi M_{\{0, 1, p_0, \infty\}}\\
\int_{S^2_1}  \Omega_I&= \int_{S^2_{\{0, p_0\}}} \Omega_I &= -2 \pi (-m_0 +m_1 - m_{p_0}+ m_\infty)&=2 \pi M_{\{0,p_0\}}\\
\int_{S^2_0}  \Omega_I&=\int_{S^2_{\{1, p_0\}}} \Omega_I &= -2\pi (m_0 -m_1 - m_{p_0}+ m_\infty)&=2 \pi M_{\{1,p_0\}}\\
\int_{S^2_{p_0}}  \Omega_I&=\int_{S^2_{\{0, 1\}}} \Omega_I &= -2 \pi (-m_0 -m_1 + m_{p_0}+ m_\infty)&= 2 \pi M_{\{0,1\}}
\end{array}
\end{equation}
\end{prop}

\begin{proof} 
We observe that the eight polar sections $\sigma^\pm_p$ are the same as in the proof of Proposition \ref{prop:exteriorintersection}.

As one passes from the exterior chamber of Proposition \ref{prop:exteriorintersection} to the adjacent interior chamber, all exterior $\C^\times$-fixed points remain stable. However, the description of the central sphere changes entirely. From this we see that the flows to the exterior $\C^\times$-fixed points remain unchanged:
 \begin{itemize}
   \item $\sigma_0^-$ all flows to the exterior $\C^\times$-fixed point with $u=0$. One can indeed check that $F_1=F_{p_0}=\langle e_2\rangle$ and $F_0=F_\infty=\langle e_1\rangle$.
    \item$\sigma_1^-$ all flows to the exterior $\C^\times$-fixed point with $u=1$. One can indeed check that $F_0=F_{p_0}=\langle e_2\rangle$ and $F_1=F_\infty=\langle e_1\rangle$.  
        \item $\sigma_{p_0}^-$ all flows to the exterior $\C^\times$-fixed point with $u=p_0$. One can indeed check that $F_0=F_1=\langle e_2\rangle$ and $F_{p_0}=F_\infty=\langle e_1\rangle$. 
        \item $\sigma_{\infty}^+$ all flows to the exterior $\C^\times$-fixed point at $u=\infty$.
        \end{itemize}
In fact, $\sigma_0^- = \tau_0^+, \sigma_1^-=\tau_1^+, \sigma_{p_0}^-=\tau_{p_0}^+, \sigma_\infty^+ =\tau_\infty^+$. 

It remains to consider the flows of these other four polar sections under the $\C^\times$-flow.  Since each point of these four polar sections flows somewhere on the central sphere, the $\C^\times$-limit is simply the underlying stable parabolic bundle with zero Higgs field. In our particular chamber, this is determined by the cross-ratio of the flags (see \ref{sec:centralsphereinint}).
From Corollary \ref{cor:flags}, we recall that for the big strata, the flags are    \begin{equation*}
    F_0 = \begin{pmatrix} \frac{-m_0 p_0 +w}{u} \\ 1 \end{pmatrix}, \quad 
      F_1 = 
  \begin{pmatrix} \frac{m_1(p_0-1) + m_\infty +w}{u-1} \\ 1 \end{pmatrix}, \quad 
        F_{p_0} = \begin{pmatrix} \frac{-m_{p_0} p_0(p_0-1) + m_\infty p_0^2 +w}{u-p_0}\\ 1 \end{pmatrix}, \quad 
          F_\infty = \begin{pmatrix} 1\\ 0 \end{pmatrix}.
   \end{equation*}
   Note that the formula for $F_p$ is not defined on $\sigma^+_p$ for $p \in \{0, 1, p_0\}$, however, we can use L'H\^{o}pital's Rule to extend the formula for $F_p$ to $\sigma^+_p$
\begin{itemize}
\item On $\sigma^+_0$, \[
F_0 = \begin{pmatrix} \frac{f'(0) + \beta p_0}{2 m_0 p_0} \\ 1.\end{pmatrix}
\]
To see this observe by L'H\^{o}pital's Rule
\begin{align*} F_0\Big|_{\sigma^+_0} &= \frac{-m_0 p_0 + w}{u}\Big|_{\sigma^+_0} \\
& =\frac{\frac{\partial}{\partial \beta}(m_0 p_0 + w)}{\frac{\partial}{\partial \beta}(u)}\Big|_{\sigma^+_0}\\
&=\frac{\partial_\beta w}{\partial_\beta u} \Big|_{\sigma^+_0}\\
&\overset{(1)}{=}\frac{f'_{\mathbf{m}}(0)  + \beta p_0}{2m_0 p_0}
\end{align*}
In (1), we use $\partial_\beta$ of the defining equation of the elliptic curve:
\[ 0=f'_{\mathbf{m}}(u) \partial_\beta u + u(u-1)(u-p_0)+\beta \partial_u \left( u(u-1)(u-p_0) \right) \partial_\beta u -2 (m_\infty u^2 + w)( 2 m_\infty u \partial_\beta u + \partial_\beta w).\]
Evaluated on $\sigma^+_0$ ($u=0$, $w=m_0 p_0$), this gives 
\[ 0 = \left( f'_{\mathbf{m}}(0)  + \beta p_0 \right)\partial_\beta u    - 2 (m_0p_0)\partial_\beta w. \]
   \item On $\sigma^+_1$,
   we compute 
   \begin{align*}
F_1\Big|_{\sigma^+_1} &= 
\frac{\partial_\beta w}{\partial_\beta (u-1)} \Big|_{\sigma^+_1}\\  
&=  \frac{f'_{\mathbf{m}}(1) + \beta(1-p_0)- 2 m_1(1-p_0) 2 m_\infty}{2m_1(1-p_0)} 
 \end{align*}
   \item On $\sigma^+_{p_0}$, we compute
      \begin{align*}
F_1\Big|_{\sigma^+_{p_0}} &= 
\frac{\partial_\beta w}{\partial_\beta (u-p_0)} \Big|_{\sigma^+_{p_0}}\\  
&=  \frac{f'_{\mathbf{m}}(p_0) + \beta p_0(p_0-1)- 2 m_{p_0} p_0(p_0-1)2 m_\infty p_0}{2 m_{p_0} p_0(p_0-1)} 
 \end{align*}
   \end{itemize}
   Now, the $\beta$ dependence of $F_{p}$ on $\sigma^+_p$ for $p \in \{0, 1, p_0\}$ is manifest. The other flags are independent of $\beta$.
   
Now, we look at the image under the $\C^\times$-limit of the polar sections $\sigma^+_0, \sigma^+_1, \sigma^+_{p_0}, \sigma_\infty^-$ , by considering the cross-ratio of the points.
Note that since $F_\infty=\infty$, the cross-ratio is 
\[ (F_\infty, F_0; F_1, F_{p_0})= \frac{F_{p_0} - F_0}{F_1-F_0}.\]
\begin{itemize}
\item Let's first consider $\sigma^+_0$, which we note is parameterized by $\beta \in \C$. Note that $F_1, F_{p_0}, F_\infty$ are all fixed on $\sigma^-_0$, so the cross-ratio of the flag only sees $\beta$ through $F_0$.
The important thing is that the cross-ratio of the flags has the shape of a conformal transformation:
\begin{equation} (F_\infty, F_0; F_1, F_{p_0}) = \frac{a_0 \beta + a_1}{a_2 \beta + a_3}\label{eq:conformal}\end{equation}
for fixed values $a_0, a_1, a_2, a_3$. 
Consequently, its image is all of $\mathbb{CP}^1$, minus the point 
\[ \lim_{\beta \to \infty} (F_\infty, F_0; F_1, F_{p_0}) = \frac{a_0}{a_2}, \]
Since $\beta$ only appears in $F_0$, we can see that 
\[ \lim_{\beta \to \infty} (F_\infty, F_0; F_1, F_{p_0}) =\frac{-F_0}{-F_0} = 1.\]
\item Now, let's consider $\sigma^+_1$. Since $\beta$ only appears in $F_1$, we see
\[ \lim_{\beta \to \infty} (F_\infty, F_0; F_1, F_{p_0}) =\frac{0}{F_1} = 0.\]
\item Now, let's consider $\sigma^+_{p_0}$. Since $\beta$ only appears in $F_{p_0}$, we see 
\[ \lim_{\beta \to \infty} (F_\infty, F_0; F_1, F_{p_0}) =\frac{F_p}{0}= \infty.\]
\item 
Lastly, we consider $\sigma_\infty^-$, the additional section in the big stratum. Note that from Corollary \ref{cor:flags}, $\beta$ appears in both $F_1$ (with coefficient $-\frac{1}{2m_\infty}$) and $F_{p_0}$ (with coefficient $-\frac{p_0}{2m_\infty}$), but not in $F_0, F_\infty$. Thus, 
\[ \lim_{\beta \to \infty} (F_\infty, F_0; F_1, F_{p_0}) = \frac{-\frac{p_0}{2m_\infty}}{-\frac{1}{2m_\infty}} = p_0. \]
\end{itemize}


We must match these flag arrangements to a value of $u$. 
We observe that when $\mathbf{m}=\mathbf{0}$, the cross-ratio of the flags is independent of $\beta$ and is simply 
\[(F_\infty, F_0; F_1, F_{p_0} ) = \frac{p_0(u-1)}{(u-p_0)}. \]
We indeed see that $u=0$ corresponds to the flag $1$; $u=1$ corresponds to the flag $0$; $u=p_0$ corresponds to the flag $\infty$; $u=\infty$ corresponds to the flag $p_0$. 
Thus, we see that we get the following intersections as the polar sections (each modeled on $\beta \in \C$) get wrapped around the central sphere minus a single point of $D$.
\begin{equation}
\begin{tabular}{  c | c | c |c|c }
    I & $\sigma^+_0$ & $\sigma^+_1$ & $\sigma^+_{p_0}$ & $\sigma^-_\infty$ \\ \hline
    $\tau_0^-$ & 0 & 1 & 1 & 1\\ \hline
    $\tau_1^-$ & 1 & 0 & 1 & 1\\ \hline
    $\tau_{p_0}^-$ & 1 & 1  & 0 & 1\\ \hline
    $\tau_\infty^-$ & 1 & 1  & 1 & 0\\
\end{tabular}
\end{equation}
The intersection numbers for $\Sigma_p = \sigma_p^+- \sigma_p^-$ are thus given as follows:

\begin{center}

\begin{tabular}{l c r}
	
\begin{tabular}{  c | c | c |c|c }
    I & $\Sigma_0$ & $\Sigma_1$ & $\Sigma_{p_0}$ & $\Sigma_\infty$ \\ \hline
    $\tau_0^+$ & -1 & 0 & 0 & 0\\ \hline
    $\tau_1^+$ & 0 & -1 & 0 & 0\\ \hline
    $\tau_{p_0}^+$ & 0 & 0  & -1 & 0\\ \hline
    $\tau_\infty^+$ & 0 & 0  & 0 & 1\\
\end{tabular}

&  and &

\begin{tabular}{  c | c | c |c|c }
    I & $\Sigma_0$ & $\Sigma_1$ & $\Sigma_{p_0}$ & $\Sigma_\infty$ \\ \hline
    $\tau_0^-$ & 0 & -1 & -1 & 1\\ \hline
    $\tau_1^-$ & -1 & 0 & -1 & 1\\ \hline
    $\tau_{p_0}^-$ & -1 & -1  & 0 & 1\\ \hline
    $\tau_\infty^-$ & -1 & -1  & -1 & 0.\\
\end{tabular}

\end{tabular}

\end{center}

Thus, the intersection numbers for $S_p = \tau_p^--\tau_p^+$ are 

\begin{center}

 \begin{tabular}{  c | c | c |c|c }
    I & $\Sigma_0$ & $\Sigma_1$ & $\Sigma_{p_0}$ & $\Sigma_\infty$ \\ \hline
    $S_0$ & 1 & -1 & -1 & 1\\ \hline
    $S_1$ & -1 & 1 & -1 & 1\\ \hline
    $S_{p_0}$ & -1 & -1  & 1 & 1\\ \hline
    $S_\infty$ & -1 & -1  & -1 & -1. \\
\end{tabular}

\end{center}

\end{proof}

\end{document}
